\documentclass[a4paper, 10pt]{article}
\usepackage{scrextend}
\changefontsizes{9pt}
% Per-subject config for this repo. See vendor/template/course-config.tex.example
% for what each macro is used for.

\newcommand{\CourseCode}{2110473}
\newcommand{\CourseName}{Fault Tolerant Computing}
\newcommand{\StudentID}{6733172621}
\newcommand{\StudentName}{Patthadon Phengpinij}
\newcommand{\DefaultCollaborators}{Claude (for\,\LaTeX\,styling and grammar checking)}
\newcommand{\AssignmentLabel}{Homework}

\usepackage{../vendor/template/CEDT-Assignment-style}

\begin{document}
\subject
\asgntitle{7}{}{Week 7}{\StudentID\ \StudentName}{\DefaultCollaborators}

% ================================================================================ %
%                                    Problem 01                                    %
% ================================================================================ %
\begin{problem}
A system which is operated from \( X \) to \( Y \) via the redundant non serial-parallel path as shown in the figure below.
Let each module (\( A \), \( B \), and \( K \)) has their reliabilities (\( R_a \), \( R_b \), and \( R_k \)) respectively
and ignore the reliabilities of the connections.

\par Find the system reliability.

\myspace

\begin{center}
	\begin{tikzpicture}[
		line width=0.9pt,
		signal/.style={
				draw=blue!55!cyan,
				line width=1pt
			},
		arrow/.style={
		signal,
		-{Latex[length=2.2mm,width=1.7mm]}
		},
		block/.style={
				draw=blue!25!black,
				line width=0.9pt,
				fill=white,
				minimum width=1.4cm,
				minimum height=0.95cm,
				font=\Large,
				inner sep=0pt
			},
		junction/.style={
				circle,
				fill=blue!25!black,
				inner sep=1.3pt
			}
		]

		% Blocks A
		\node[block] (A1) at (1.65,3.6) {A};
		\node[block] (A2) at (1.65,1.8) {A};
		\node[block] (A3) at (1.65,0.0) {A};

		% Blocks B
		\node[block] (B1) at (6.55,3.6) {B};
		\node[block] (B2) at (6.55,1.8) {B};
		\node[block] (B3) at (6.55,0.0) {B};

		% Blocks K
		\node[block] (K1) at (4.10,2.7) {K};
		\node[block] (K2) at (4.10,0.9) {K};

		% Input X
		\node[left] at (-1.0,1.8) {\Large X};
		\draw[signal] (-1.0,1.8) -- (0,1.8);

		% Input branching
		\draw[signal] (0,3.6) -- (0,0);
		\draw[arrow] (0,3.6) -- (A1.west);
		\draw[arrow] (0,1.8) -- (A2.west);
		\draw[arrow] (0,0) -- (A3.west);

		% A to B connections
		\draw[arrow] (A1.east) -- (B1.west);
		\draw[arrow] (A2.east) -- (B2.west);
		\draw[arrow] (A3.east) -- (B3.west);

		% K1: top and middle connections
		\draw[arrow] (4.10,3.6) -- (K1.north);
		\draw[arrow] (4.10,1.8) -- (K1.south);

		% K2: middle and bottom connections
		\draw[arrow] (4.10,1.8) -- (K2.north);
		\draw[arrow] (4.10,0) -- (K2.south);

		% Output merging
		\draw[signal] (B1.east) -- (8.2,3.6);
		\draw[signal] (B2.east) -- (8.2,1.8);
		\draw[signal] (B3.east) -- (8.2,0);
		\draw[signal] (8.2,3.6) -- (8.2,0);

		% Output Y
		\draw[arrow] (8.2,1.8) -- (9.25,1.8);
		\node[right] at (9.25,1.8) {\Large Y};

		% Junction dots
		\node[junction] at (0,1.8) {};
		\node[junction] at (4.10,3.6) {};
		\node[junction] at (4.10,1.8) {};
		\node[junction] at (4.10,0) {};
		\node[junction] at (8.2,1.8) {};

	\end{tikzpicture}
\end{center}
\end{problem}

\begin{solution}
	Let $m_1,m_2,m_3$ be the midpoints of the three $A \to B$ paths (top to
	bottom). $K_1$ bridges $m_1$--$m_2$ and $K_2$ bridges $m_2$--$m_3$, so
	the network is not series-parallel. Use the decomposition (conditional
	probability) method on the $K$ modules: $X$ reaches $Y$ iff some working
	$A_i$ and some working $B_j$ have $m_i$, $m_j$ joined by working $K$'s.
	Define
	\[
		P = R_aR_b, \quad
		Q_2 = \sqbracket{1-(1-R_a)^2}\sqbracket{1-(1-R_b)^2}, \quad
		Q_3 = \sqbracket{1-(1-R_a)^3}\sqbracket{1-(1-R_b)^3}.
	\]

	\begin{itemize}
		\item \textbf{$K_1$, $K_2$ both good} ($R_k^2$): $m_1,m_2,m_3$ merge into
		      one node, giving 3 parallel $A$'s in series with 3 parallel
		      $B$'s: $R = Q_3$.
		\item \textbf{Exactly one $K$ good} ($2R_k(1-R_k)$, symmetric): two
		      paths merge into a ``2 $A$ $\parallel$ -- 2 $B$ $\parallel$''
		      block, in parallel with the remaining single $A$--$B$ path:
		      $R = 1-(1-Q_2)(1-P)$.
		\item \textbf{Both $K$ failed} ($(1-R_k)^2$): three independent
		      $A$--$B$ paths in parallel: $R = 1-(1-P)^3$.
	\end{itemize}

	By the total probability theorem,
	\[
		\boxed{
			R_{sys} = R_k^2\,Q_3
			+ 2R_k(1-R_k)\sqbracket{1-(1-Q_2)(1-R_aR_b)}
			+ (1-R_k)^2\sqbracket{1-(1-R_aR_b)^3}
		}
	\]
	\textit{Check:} $R_k=1$ gives $Q_3$ (fully cross-connected) and $R_k=0$
	gives three plain parallel paths, as expected.
\end{solution}
% ================================================================================ %


\end{document}
